% =====================================================================
\section{Materials}\label{sec:data}
% =====================================================================
\subsection{Datasets}
\paragraph{IDRiD (target dataset).}
The Indian Diabetic Retinopathy Image Dataset \citep{porwal2018idrid,porwal2020idrid} was collected at an eye clinic
in Nanded, India, with a Kowa VX-10$\alpha$ fundus camera at a resolution of $4288\times2848$ pixels and a
$50^\circ$ field of view centred near the macula. Each image is graded on the five-level ICDR scale
\citep{wilkinson2003proposed} by medical experts; a subset also has pixel-level lesion masks, which we do not use
for training. We use the publicly distributed Kaggle copy of the grading subset. Image files are matched to label
rows by file name, with a fallback that compares the digit sequence in the name; after matching and de-duplication,
\res{nidrid} graded images remain.

\paragraph{APTOS-2019 (pretraining dataset).}
The APTOS-2019 Blindness Detection dataset \citep{aptos2019} contains \res{naptos} labelled fundus photographs
collected in rural India by the Aravind Eye Hospital, graded on the same 0--4 scale. Images come from several
cameras and vary widely in size, illumination, focus and field of view, which makes the dataset a useful source of
appearance diversity. It is used only for stage-1 pretraining and never for evaluation.

Table~\ref{tab:data} summarises the grade distributions (shown graphically in Supplementary Fig.~\ref{fig:dist}). Both datasets
are imbalanced: grade 0 and grade 2 dominate, while grades 1 and 3 are comparatively rare. This imbalance motivates
the class-weighted loss (Section~\ref{sec:method}) and the use of balanced accuracy, macro-F1 and QWK alongside plain
accuracy (Section~\ref{sec:setup}).

\begin{table}[!htbp]
\centering\small
\caption{Grade distribution of the datasets used. The IDRiD test split is a stratified 15\% hold-out that is never
used for training, model selection or threshold fitting.}
\label{tab:data}
\begin{tabular}{@{}lcccccc@{}}
\toprule
Dataset & Grade 0 & Grade 1 & Grade 2 & Grade 3 & Grade 4 & Total\\
\midrule
APTOS-2019 (pretraining) & \res{nap0} & \res{nap1} & \res{nap2} & \res{nap3} & \res{nap4} & \res{naptos}\\
IDRiD (all) & \res{nid0} & \res{nid1} & \res{nid2} & \res{nid3} & \res{nid4} & \res{nidrid}\\
IDRiD test (held out) & \res{nte0} & \res{nte1} & \res{nte2} & \res{nte3} & \res{nte4} & \res{ntest}\\
\bottomrule
\end{tabular}
\end{table}



\subsection{Data partitioning}
Figure~\ref{fig:split} summarises the protocol. Before any model is trained, a stratified 15\% of IDRiD
(\res{ntest} images, random seed 42) is set aside as the test set. It is used exactly once, after all modelling
decisions have been fixed: it plays no role in training, early stopping, checkpoint selection, threshold fitting,
choice of the decision rule or any hyper-parameter. The remaining \res{ntrainval} images are divided into five
stratified folds. In each fold, four parts are used for training and one for validation (early stopping and
checkpoint selection); the validation predictions of all five folds together form the out-of-fold (OOF) set on
which the decision rule and thresholds are fitted. For APTOS-2019, a separate stratified 90/10 split is used for
checkpoint selection during pretraining.

\begin{figure}[!htbp]
\centering
\begin{tikzpicture}[node distance=4mm and 6mm]
  \node[blk,minimum width=3.0cm] (ap) {APTOS-2019\\90\% train / 10\% val};
  \node[bb,right=of ap,minimum width=3.0cm] (s1) {Stage 1\\pretraining};
  \node[blk,below=8mm of ap,minimum width=3.0cm] (id) {IDRiD (\res{nidrid})};
  \node[blk,right=of id,minimum width=3.0cm] (tv) {train+val 85\%\\5 stratified folds};
  \node[ex,below=5mm of tv,minimum width=3.0cm] (te) {test 15\%\\locked away};
  \node[bb,right=of tv,minimum width=3.0cm] (s2) {Stage 2: fine-tune\\5 fold models};
  \node[blk,right=of s2,minimum width=2.6cm] (oof) {OOF predictions\\$\rightarrow$ thresholds};
  \node[ex,below=5mm of oof,minimum width=2.6cm] (ev) {single final\\evaluation};
  \draw[arr] (ap)--(s1); \draw[arr] (id)--(tv); \draw[arr] (id.south) |- (te);
  \draw[arr] (tv)--(s2); \draw[arr] (s1.east) -| (s2.north); \draw[arr] (s2)--(oof);
  \draw[arr] (oof)--(ev); \draw[arr] (te)--(ev);
\end{tikzpicture}
\caption{Data protocol. The IDRiD test split (gold) is separated before training and enters only the final
evaluation; all data-dependent decisions use APTOS or the out-of-fold IDRiD predictions.}
\label{fig:split}
\end{figure}

\subsection{Pre-processing}
Fundus photographs from different cameras differ in resolution, field of view, black borders and illumination. Each
image is processed identically for both datasets (Fig.~\ref{fig:prep}):
\begin{enumerate}
\item images whose longer side exceeds 1600 pixels are first down-scaled with area interpolation, which keeps the
  remaining steps fast without affecting the final $448$-pixel resolution;
\item the black border is removed by cropping to the bounding box of pixels whose grey level exceeds 10;
\item the image is zero-padded to a square so that the circular fundus is not distorted, and resized to
  $448\times448$ pixels;
\item Ben Graham's local-contrast normalisation \citep{graham2015kaggle} is applied,
\begin{equation}
I_{\text{ben}} = 4\,I - 4\,(G_{\sigma} * I) + 128,\qquad \sigma = 448/30,
\label{eq:ben}
\end{equation}
  where $G_\sigma *$ denotes Gaussian blurring; subtracting the local mean removes slow illumination gradients and
  camera-specific colour casts, while the factor 4 amplifies local contrast such as small red lesions and exudates;
\item pixels outside a centred circle of radius $0.48\times448$ are set to grey (128) to suppress the bright ring
  artefact that Eq.~\eqref{eq:ben} creates at the edge of the fundus.
\end{enumerate}
As an alternative pre-processing, evaluated in the ablation study, we apply contrast-limited adaptive histogram
equalisation (CLAHE) \citep{zuiderveld1994contrast} to the green channel (clip limit 2, $8\times8$ tiles), where
vessels and red lesions have the highest contrast, followed by an edge-preserving bilateral filter.
Pre-processed images are cached once and then normalised with ImageNet channel statistics.

\begin{figure}[!t]
\centering
\fig[0.26]{samples_raw.png}{0.8\linewidth}{Raw IDRiD images of grades 0, 2, 3, 4.}\\[2pt]
\fig[0.26]{samples_pre.png}{0.8\linewidth}{The same images after pre-processing.}
\caption{Example IDRiD images before (top) and after (bottom) cropping, padding and Ben Graham normalisation.}
\label{fig:prep}
\end{figure}

\subsection{Data augmentation}
To reduce over-fitting on a few hundred images, training images are augmented on the fly with random horizontal
and vertical flips; random rotation in $[-180^\circ,180^\circ]$, which is valid because the circular fundus has no
canonical orientation once the image is padded; translation of up to 3\%; scaling by a factor in $[0.95,1.10]$;
colour jitter (brightness and contrast 0.2, saturation 0.1, hue 0.02); and random erasing \citep{zhong2020random}
of 2--8\% of the image with probability 0.2, which discourages reliance on a single region. Validation and test
images are not augmented, except for the deterministic flips used for test-time augmentation
(Section~\ref{sec:method}).
\sectionend{data}
